%%%%%%%%%%%%%%%%%%%%%%%%%%%%%%%%%%%%%%%%%%%%%%%
%
% Template per Elaborato di Laurea
% DISI - Dipartimento di Ingegneria e Scienza dell’Informazione
%
% update 2015-09-10
%
% Per la generazione corretta del 
% pdflatex nome_file.tex
% bibtex nome_file.aux
% pdflatex nome_file.tex
% pdflatex nome_file.tex
%
%%%%%%%%%%%%%%%%%%%%%%%%%%%%%%%%%%%%%%%%%%%%%%%

% formato FRONTE RETRO
\documentclass[epsfig,a4paper,11pt,titlepage,twoside,openany]{book}
\usepackage{epsfig}
\usepackage{plain}
\usepackage{setspace}
\usepackage[paperheight=29.7cm,paperwidth=21cm,outer=1.5cm,inner=2.5cm,top=2cm,bottom=2cm]{geometry} % per definizione layout
\usepackage{titlesec} % per formato custom dei titoli dei capitoli

%%%%%%%%%%%%%%
% supporto lettere accentate
%
%\usepackage[latin1]{inputenc} % per Windows;
\usepackage[utf8x]{inputenc} % per Linux (richiede il pacchetto unicode);
%\usepackage[applemac]{inputenc} % per Mac.

\singlespacing

\usepackage[italian]{babel}

\begin{document}

  % nessuna numerazione
  \pagenumbering{gobble} 
  \pagestyle{plain}

\thispagestyle{empty}

\begin{center}
  \begin{figure}[h!]
    \centerline{\psfig{file=marchio_unitrento_colore_it_202002.eps,width=0.6\textwidth}}
  \end{figure}

  \vspace{2 cm} 

  \LARGE{Department of Information Engineering and Computer Science\\}

  \vspace{1 cm} 
  \Large{Bachelor's Degree in\\
    Computer Science
  }

  \vspace{2 cm} 
  \Large\textsc{Final dissertation\\} 
  \vspace{1 cm} 
  \Huge\textsc{Drone Surveillance: a ROS2 solution for ground robots autonomous navigation\\}
  \Large{\it{Both for real and simulated environments (?)}}


  \vspace{2 cm} 
  \begin{tabular*}{\textwidth}{ c @{\extracolsep{\fill}} c }
  \Large{Supervisor} & \Large{Student}\\
  \Large{Palopoli Luigi}& \Large{Franzin Mattia}\\
  \end{tabular*}

  \vspace{2 cm} 

  \Large{Academic year 2021/2022}
  
\end{center}



  \clearpage
 
%%%%%%%%%%%%%%%%%%%%%%%%%%%%%%%%%%%%%%%%%%%%%%%%%%%%%%%%%%%%%%%%%%%%%%%%%%
%%%%%%%%%%%%%%%%%%%%%%%%%%%%%%%%%%%%%%%%%%%%%%%%%%%%%%%%%%%%%%%%%%%%%%%%%%
%% Nota
%%%%%%%%%%%%%%%%%%%%%%%%%%%%%%%%%%%%%%%%%%%%%%%%%%%%%%%%%%%%%%%%%%%%%%%%%%
%% Sezione Ringraziamenti opzionale
%%%%%%%%%%%%%%%%%%%%%%%%%%%%%%%%%%%%%%%%%%%%%%%%%%%%%%%%%%%%%%%%%%%%%%%%%%
%%%%%%%%%%%%%%%%%%%%%%%%%%%%%%%%%%%%%%%%%%%%%%%%%%%%%%%%%%%%%%%%%%%%%%%%%%
  \thispagestyle{empty}

\begin{center}
  {\bf \Huge Acknowledgements}
\end{center}

\vspace{4cm}


\emph{
  ...thanks to...
}

  \clearpage
  \pagestyle{plain} % nessuna intestazione e pie pagina con numero al centro

  
  % inizio numerazione pagine in numeri arabi
  \mainmatter

%%%%%%%%%%%%%%%%%%%%%%%%%%%%%%%%%%%%%%%%%%%%%%%%%%%%%%%%%%%%%%%%%%%%%%%%%%
%%%%%%%%%%%%%%%%%%%%%%%%%%%%%%%%%%%%%%%%%%%%%%%%%%%%%%%%%%%%%%%%%%%%%%%%%%
%% Nota
%%%%%%%%%%%%%%%%%%%%%%%%%%%%%%%%%%%%%%%%%%%%%%%%%%%%%%%%%%%%%%%%%%%%%%%%%%
%% Si ricorda che il numero massimo di facciate e' 30.
%% Nel conteggio delle facciate sono incluse 
%%   indice
%%   sommario
%%   capitoli
%% Dal conteggio delle facciate sono escluse
%%   frontespizio
%%   ringraziamenti
%%   allegati    
%%%%%%%%%%%%%%%%%%%%%%%%%%%%%%%%%%%%%%%%%%%%%%%%%%%%%%%%%%%%%%%%%%%%%%%%%%
%%%%%%%%%%%%%%%%%%%%%%%%%%%%%%%%%%%%%%%%%%%%%%%%%%%%%%%%%%%%%%%%%%%%%%%%%%

    % indice
    \tableofcontents
    \clearpage
    
    
          
    % gruppo per definizone di successione capitoli senza interruzione di pagina
    \begingroup
      % nessuna interruzione di pagina tra capitoli
      % ridefinizione dei comandi di clear page
      \renewcommand{\cleardoublepage}{} 
      \renewcommand{\clearpage}{} 
      % redefinizione del formato del titolo del capitolo
      % da formato
      %   Capitolo X
      %   Titolo capitolo
      % a formato
      %   X   Titolo capitolo
      
      \titleformat{\chapter}
        {\normalfont\Huge\bfseries}{\thechapter}{1em}{}
        
      \titlespacing*{\chapter}{0pt}{0.59in}{0.02in}
      \titlespacing*{\section}{0pt}{0.20in}{0.02in}
      \titlespacing*{\subsection}{0pt}{0.10in}{0.02in}
      
      % sommario
      \chapter*{Summary} % senza numerazione
\label{sommario}

\addcontentsline{toc}{chapter}{Summary} % da aggiungere comunque all'indice

\section*{Introduction}

The purpose of this project is to develop (developing?) a fullstack (?) system for surveillance proposes using mobile robots, both ground and flying ones.
Since the project is quite big (?), it is shared with other two students with each one of us focusing into a specific subtask (?). It is constituted of: 

\begin{itemize}
  \item a planning system, responsible for dispatching the surveillance tasks to UGVs and/or UAVs in the best possible way based on their battery charge, giving them some waypoints to follow in order to follow the optimal path.
  \item a solution to autonomously control the drones (UAVs robots)
  \item a solution to autonomously control ground robots avoiding dynamic obstacles (UGVs robots)
\end{itemize}

Basically, all the work is done using ROS2\footnote{The Robot Operating System (ROS) is a set of software libraries and tools for building robot applications (https://docs.ros.org/en/foxy/index.html, 27/07)} which gives us the possibility to work separately on your own project, and since everyone creates his custom packages, at the end we only need to put everything together.

A more detailed description can be found in the first chapter.

The reason behind my choice about the internship and this consequent thesis is leaded by my growing interest in this topic: just some month before taking part in this work I attended a course about robotic fundamentals\footnote{Fondamenti di robotica (https://www.esse3.unitn.it/Guide/PaginaADErogata.do?ad\_er\_id=2021*N0*N0*S1*51102 *94282\&ANNO\_ACCADEMICO=2021\&mostra\_percorsi=S, 27/07)} in which I had to work in a custom project similar to this, but using a manipulator.

For me, is the possibility to give to some inanimate object, like a wheeled robot or a manipulator, something that could be described as intelligence, the ability to perform some tasks in response of other ones, figuring out which ones are the best for every particular situation. Another aspect I really like is the (inbuilt) utility which carries (within) itself: this is only an educational project, but it is easy to imagine an industrial application, for example for patrol purposes. // ?

Some challenges I decided to took are the choice of the new version of ROS, with some great improves respect to the older one, but with less community support, and working for the first time with a mobile robot.

Sommario è un breve riassunto del lavoro svolto dove si descrive l'obiettivo, l'oggetto della tesi, le 
metodologie e le tecniche usate, i dati elaborati e la spiegazione delle conclusioni alle quali siete arrivati.  

% Il sommario dell’elaborato consiste al massimo di 3 pagine e deve contenere le seguenti informazioni:
% \begin{itemize}
%   \item contesto e motivazioni 
%   \item breve riassunto del problema affrontato
%   \item tecniche utilizzate e/o sviluppate
%   \item risultati raggiunti, sottolineando il contributo personale del laureando/a
% \end{itemize}

% Tipo:
% - descrizione tesi e interesse robotica
% - far muovere robot in un ambiente seguendo le istruzioni ricevute ed evitando ostacoli non previsti
% - ROS2, nav2, waypoint follower
% - il robot si muove e fa foto

\section*{Description of other students subtasks}

What follows is a brief description of the work done by the other two students to have a better and clear idea of the workflow of the entire project.

\subsection*{Planning system}

To be done...

\subsection*{Drones control}

To be done...

%%%%%%%%%%%%%%%%%%%%%%%%%%%%%%%%%%%%%%%%%%%%%%%%%%%%%%%%%%%%%%%%%%%%%%%%%%
%%%%%%%%%%%%%%%%%%%%%%%%%%%%%%%%%%%%%%%%%%%%%%%%%%%%%%%%%%%%%%%%%%%%%%%%%%
%% Nota
%%%%%%%%%%%%%%%%%%%%%%%%%%%%%%%%%%%%%%%%%%%%%%%%%%%%%%%%%%%%%%%%%%%%%%%%%%
%% Sommario e' un breve riassunto del lavoro svolto dove si descrive 
%% l’obiettivo, l’oggetto della tesi, le metodologie e 
%% le tecniche usate, i dati elaborati e la spiegazione delle conclusioni 
%% alle quali siete arrivati.
%% Il sommario dell’elaborato consiste al massimo di 3 pagine e deve contenere le seguenti informazioni: 
%%   contesto e motivazioni
%%   breve riassunto del problema affrontato
%%   tecniche utilizzate e/o sviluppate
%%   risultati raggiunti, sottolineando il contributo personale del laureando/a
%%%%%%%%%%%%%%%%%%%%%%%%%%%%%%%%%%%%%%%%%%%%%%%%%%%%%%%%%%%%%%%%%%%%%%%%%%
%%%%%%%%%%%%%%%%%%%%%%%%%%%%%%%%%%%%%%%%%%%%%%%%%%%%%%%%%%%%%%%%%%%%%%%%%%      
      
      %%%%%%%%%%%%%%%%%%%%%%%%%%%%%%%%
      % lista dei capitoli
      %
      % \input oppure \include
      %
      \chapter{In ante nulla, vestibulum a}
\label{cha:intro}

Lorem ipsum dolor sit amet, consectetur adipiscing elit. Donec sed nunc orci. Aliquam nec nisl vitae sapien pulvinar dictum quis non urna. Suspendisse at dui a erat aliquam vestibulum. Quisque ultrices pellentesque pellentesque. Pellentesque egestas quam sed blandit tempus. Sed congue nec risus posuere euismod. Maecenas ut lacus id mauris sagittis egestas a eu dui. Class aptent taciti sociosqu ad litora torquent per conubia nostra, per inceptos himenaeos. Pellentesque at ultrices tellus. Ut eu purus eget sem iaculis ultricies sed non lorem. Curabitur gravida dui eget ex vestibulum venenatis. Phasellus gravida tellus velit, non eleifend justo lobortis eget. 
\cite{coulouris}

Donec eu ipsum id lorem consectetur luctus ac a nisi. Curabitur volutpat, metus id porta ultrices, felis lacus consectetur justo, ut gravida arcu ex in purus. Pellentesque vitae sapien ac nisl porttitor pellentesque eu sed elit. Sed maximus lectus eu eros ultricies accumsan. Quisque congue, nisi in dictum cursus, ante nisl molestie eros, in ultrices eros tellus sit amet augue. Interdum et malesuada fames ac ante ipsum primis in faucibus. Nam finibus leo sit amet purus vehicula, eget facilisis turpis convallis. Vivamus varius tincidunt turpis, id venenatis arcu maximus ut. Aenean euismod eros ac nibh facilisis, nec imperdiet ex suscipit.
\cite{dalal}


\section{Pellentesque habitant morbi tristique senectus}
\label{sec:context}

Lorem ipsum dolor sit amet, consectetur adipiscing elit. Donec sed nunc orci. Aliquam nec nisl vitae sapien pulvinar dictum quis non urna. Suspendisse at dui a erat aliquam vestibulum. Quisque ultrices pellentesque pellentesque. Pellentesque egestas quam sed blandit tempus. Sed congue nec risus posuere euismod. Maecenas ut lacus id mauris sagittis egestas a eu dui. Class aptent taciti sociosqu ad litora torquent per conubia nostra, per inceptos himenaeos. Pellentesque at ultrices tellus. Ut eu purus eget sem iaculis ultricies sed non lorem. Curabitur gravida dui eget ex vestibulum venenatis. Phasellus gravida tellus velit, non eleifend justo lobortis eget.
\cite{ictbusiness}
\cite{donoho}

\section{Nullam et justo vitae nisi}
\label{sec:problem}

Lorem ipsum dolor sit amet, consectetur adipiscing elit. Donec sed nunc orci. Aliquam nec nisl vitae sapien pulvinar dictum quis non urna. Suspendisse at dui a erat aliquam vestibulum. Quisque ultrices pellentesque pellentesque. Pellentesque egestas quam sed blandit tempus. Sed congue nec risus posuere euismod. Maecenas ut lacus id mauris sagittis egestas a eu dui. Class aptent taciti sociosqu ad litora torquent per conubia nostra, per inceptos himenaeos. Pellentesque at ultrices tellus. Ut eu purus eget sem iaculis ultricies sed non lorem. Curabitur gravida dui eget ex vestibulum venenatis. Phasellus gravida tellus velit, non eleifend justo lobortis eget.



      \chapter{Proin rhoncus a sapien in.}
\label{cha:789}
Lorem ipsum dolor sit amet, consectetur adipiscing elit. Donec sed nunc orci. Aliquam nec nisl vitae sapien pulvinar dictum quis non urna. Suspendisse at dui a erat aliquam vestibulum. Quisque ultrices pellentesque pellentesque. Pellentesque egestas quam sed blandit tempus. Sed congue nec risus posuere euismod. Maecenas ut lacus id mauris sagittis egestas a eu dui. Class aptent taciti sociosqu ad litora torquent per conubia nostra, per inceptos himenaeos. Pellentesque at ultrices tellus. Ut eu purus eget sem iaculis ultricies sed non lorem. Curabitur gravida dui eget ex vestibulum venenatis. Phasellus gravida tellus velit, non eleifend justo lobortis eget. 


\section{Cras in aliquam quam, et}
\label{sec:456}
Lorem ipsum dolor sit amet, consectetur adipiscing elit. Donec sed nunc orci. Aliquam nec nisl vitae sapien pulvinar dictum quis non urna. Suspendisse at dui a erat aliquam vestibulum. Quisque ultrices pellentesque pellentesque. Pellentesque egestas quam sed blandit tempus. Sed congue nec risus posuere euismod. Maecenas ut lacus id mauris sagittis egestas a eu dui. Class aptent taciti sociosqu ad litora torquent per conubia nostra, per inceptos himenaeos. Pellentesque at ultrices tellus. Ut eu purus eget sem iaculis ultricies sed non lorem. Curabitur gravida dui eget ex vestibulum venenatis. Phasellus gravida tellus velit, non eleifend justo lobortis eget.


\subsection{Sed pulvinar placerat enim, a}
\label{sec:00456}
Lorem ipsum dolor sit amet, consectetur adipiscing elit. Donec sed nunc orci. Aliquam nec nisl vitae sapien pulvinar dictum quis non urna. Suspendisse at dui a erat aliquam vestibulum. Quisque ultrices pellentesque pellentesque. Pellentesque egestas quam sed blandit tempus. Sed congue nec risus posuere euismod. Maecenas ut lacus id mauris sagittis egestas a eu dui. Class aptent taciti sociosqu ad litora torquent per conubia nostra, per inceptos himenaeos. Pellentesque at ultrices tellus. Ut eu purus eget sem iaculis ultricies sed non lorem. Curabitur gravida dui eget ex vestibulum venenatis. Phasellus gravida tellus velit, non eleifend justo lobortis eget.


\section{Vivamus hendrerit imperdiet ex. Vivamus}
\label{sec:123}
Lorem ipsum dolor sit amet, consectetur adipiscing elit. Donec sed nunc orci. Aliquam nec nisl vitae sapien pulvinar dictum quis non urna. Suspendisse at dui a erat aliquam vestibulum. Quisque ultrices pellentesque pellentesque. Pellentesque egestas quam sed blandit tempus. Sed congue nec risus posuere euismod. Maecenas ut lacus id mauris sagittis egestas a eu dui. Class aptent taciti sociosqu ad litora torquent per conubia nostra, per inceptos himenaeos. Pellentesque at ultrices tellus. Ut eu purus eget sem iaculis ultricies sed non lorem. Curabitur gravida dui eget ex vestibulum venenatis. Phasellus gravida tellus velit, non eleifend justo lobortis eget.



      \chapter{Conclusioni}
\label{cha:conclusioni}
Lorem ipsum dolor sit amet, consectetur adipiscing elit. Donec sed nunc orci. Aliquam nec nisl vitae sapien pulvinar dictum quis non urna. Suspendisse at dui a erat aliquam vestibulum. Quisque ultrices pellentesque pellentesque. Pellentesque egestas quam sed blandit tempus. Sed congue nec risus posuere euismod. Maecenas ut lacus id mauris sagittis egestas a eu dui. Class aptent taciti sociosqu ad litora torquent per conubia nostra, per inceptos himenaeos. Pellentesque at ultrices tellus. Ut eu purus eget sem iaculis ultricies sed non lorem. Curabitur gravida dui eget ex vestibulum venenatis. Phasellus gravida tellus velit, non eleifend justo lobortis eget. 


      %\input{capitolo4}
      
      
    \endgroup


    % bibliografia in formato bibtex
    %
    % aggiunta del capitolo nell'indice
    \addcontentsline{toc}{chapter}{Bibliografia}
    % stile con ordinamento alfabetico in funzione degli autori
    \bibliographystyle{plain}
    \bibliography{biblio}
%%%%%%%%%%%%%%%%%%%%%%%%%%%%%%%%%%%%%%%%%%%%%%%%%%%%%%%%%%%%%%%%%%%%%%%%%%
%%%%%%%%%%%%%%%%%%%%%%%%%%%%%%%%%%%%%%%%%%%%%%%%%%%%%%%%%%%%%%%%%%%%%%%%%%
%% Nota
%%%%%%%%%%%%%%%%%%%%%%%%%%%%%%%%%%%%%%%%%%%%%%%%%%%%%%%%%%%%%%%%%%%%%%%%%%
%% Nella bibliografia devono essere riportati tutte le fonti consultate 
%% per lo svolgimento della tesi. La bibliografia deve essere redatta 
%% in ordine alfabetico sul cognome del primo autore. 
%% 
%% La forma della citazione bibliografica va inserita secondo la fonte utilizzata:
%% 
%% LIBRI
%% Cognome e iniziale del nome autore/autori, la data di edizione, titolo, casa editrice, eventuale numero dell’edizione. 
%% 
%% ARTICOLI DI RIVISTA
%% Cognome e iniziale del nome autore/autori, titolo articolo, titolo rivista, volume, numero, numero di pagine.
%% 
%% ARTICOLI DI CONFERENZA
%% Cognome e iniziale del nome autore/autori (anno), titolo articolo, titolo conferenza, luogo della conferenza (città e paese), date della conferenza, numero di pagine. 
%% 
%% SITOGRAFIA
%% La sitografia contiene un elenco di indirizzi Web consultati e disposti in ordine alfabetico. 
%% E’ necessario:
%%   Copiare la URL (l’indirizzo web) specifica della pagina consultata
%%   Se disponibile, indicare il cognome e nome dell’autore, il titolo ed eventuale sottotitolo del testo
%%   Se disponibile, inserire la data di ultima consultazione della risorsa (gg/mm/aaaa).    
%%%%%%%%%%%%%%%%%%%%%%%%%%%%%%%%%%%%%%%%%%%%%%%%%%%%%%%%%%%%%%%%%%%%%%%%%%
%%%%%%%%%%%%%%%%%%%%%%%%%%%%%%%%%%%%%%%%%%%%%%%%%%%%%%%%%%%%%%%%%%%%%%%%%%
    

    \titleformat{\chapter}
        {\normalfont\Huge\bfseries}{Allegato \thechapter}{1em}{}
    % sezione Allegati - opzionale
    \appendix
    \chapter{Docker Image}
\label{cha:dockerfile}

The project depends on the following packages:

\begin{itemize}
    \item todo
\end{itemize}

If you prefer you could use the following Dockerfile to create an image starting from Ubuntu 20.04 with the required packages:

\begin{lstlisting}[
  language=Dockerfile,
  style=DockerStyle,
  caption={Dockerfile used}
  ]
FROM osrf/ros:foxy-ros1-bridge

ENV DEBIAN_FRONTEND noninteractive

# Fix expired keys

RUN apt-key adv --keyserver keyserver.ubuntu.com --recv-keys F42ED6FBAB17C654


RUN \
  apt-get update && \
  apt-get -y install apt-utils

RUN \
  apt-get -y install nano python3-pip libzmq3-dev

# ROS1 dependencies for shelfino
  
RUN \
  apt-get -y install ros-noetic-tf2-bullet \
                     ros-noetic-grid-map-ros \
                     ros-noetic-voxel-grid \
                     ros-noetic-amcl \
                     ros-noetic-costmap-2d \
                     ros-noetic-map-server \
                     ros-noetic-depthimage-to-laserscan \
                     ros-noetic-slam-toolbox --fix-missing

# ROS2 dependencies for current project

RUN \
  apt-get install -y python3-colcon-common-extensions \
                     ros-foxy-test-msgs \
                     ros-foxy-gazebo-ros-pkgs \
                     ros-foxy-xacro \
                     ros-foxy-joint-state-publisher-gui \
                     ros-foxy-rqt-robot-steering \
                     ros-foxy-nav2-bringup \
                     ros-foxy-robot-localization \
                     ros-foxy-lifecycle-msgs \
                     ros-foxy-rviz2

# ROS2 dependencies for planning

RUN apt-get install -y \
  python3-colcon-common-extensions \
  ros-foxy-test-msgs \
  ros-foxy-rosbridge-server \
  ros-foxy-plansys2-*

RUN apt-get install -y gdb # and more

RUN sed -i 's/(ALL:ALL) ALL/(ALL) NOPASSWD: ALL/' /etc/sudoers # Enable sudo without password

# Add user with home folder
RUN useradd -ms /bin/bash -G sudo ros2
USER ros2
WORKDIR /home/ros2

RUN echo 'alias noetic="source /opt/ros/noetic/setup.bash"' >> ~/.bashrc
RUN echo 'alias foxy="source /opt/ros/foxy/setup.bash"' >> ~/.bashrc
RUN echo 'alias rosplz1="source ./devel/setup.bash"' >> ~/.bashrc
RUN echo 'alias rosplz2="source ./install/setup.bash"' >> ~/.bashrc

RUN echo 'alias rosinstall="rosdep install -y -r -q --from-paths src \
                            --ignore-src --rosdistro foxy"' >> ~/.bashrc
RUN echo "source /usr/share/colcon_cd/function/colcon_cd.sh" >> ~/.bashrc
RUN echo "export GAZEBO_MODEL_PATH=~/workspace2/src/g_robot/models" >> ~/.bashrc

ENTRYPOINT ["bash"]
\end{lstlisting}

\chapter{Planning Bridge deepening}
\label{cha:deepening}

\newsavebox\poseserver

\noindent\begin{lrbox}{\poseserver}
  \noindent\begin{lstlisting}[language=C, style=CStyleTextWidth]
while (rclcpp::ok() && !listener_.buffer_.canTransform("map", "base_link", tf2::TimePoint())) rate.sleep();

geometry_msgs::msg::TransformStamped transform_ = listener_.buffer_.lookupTransform("map", "base_link", tf2::TimePoint());

response->x = transform_.transform.translation.x;
response->y = transform_.transform.translation.y;
  \end{lstlisting}
\end{lrbox}

\newsavebox\navigationclient

\noindent\begin{lrbox}{\navigationclient}
  \noindent\begin{lstlisting}[language=C, style=CStyleNoFrame]
goal_->pose.position.x = request->x;
goal_->pose.position.y = request->y;
goal_->pose.position.z = 0;
goal_->pose.orientation.x = 0.0;
goal_->pose.orientation.y = 0.0;
goal_->pose.orientation.z = 0.0;
goal_->pose.orientation.w = 1.0;
goal_->header.frame_id = "map";
goal_->header.stamp = now();
send_goal();
response->success = true;
  \end{lstlisting}
\end{lrbox}

\noindent\begin{figure}[h]
  \centering
    \noindent\begin{tikzpicture}
      \node [poseserver] (pose-server) {\usebox\poseserver};
      \node[fill=white] at (pose-server.south) {POSE SERVER};
      
      \node [inception, above left=5.7cm and -4.625cm of pose-server] (do-work) {DO WORK};
      \node [inception, above left=3.8cm and -4.625cm of pose-server] (get-pose) {GET POSE};
      \node [inception, above left=1.9cm and -4.625cm of pose-server] (continue-work) {CONTINUE WORK};
      \node[inception, inner xsep=1em, inner ysep=1em, fit=(do-work) (continue-work), above left=1.5cm and -5cm of pose-server] (plan-clients) {};
      % \node [poseserver, above=1cm of pose-server, fit=(do-work) (continue-work)] (plan-clients) {PLANNING CLIENTS};  
      \node[fill=white] at (plan-clients.south) {PLANNING CLIENTS};

      \node [navigationclient, right=5cm of plan-clients] (nav-client) {\usebox\navigationclient};
      \node[fill=white] at (nav-client.south) {NAVIGATION CLIENT};

      \node [block, above=3cm of plan-clients] (plan-server) {PLANNING SERVER};
      \node [block, right=7.9cm of plan-server] (nav-server) {NAVIGATION SERVER};

      \draw [arrow] (plan-server) -- node [left] {plansys2\_msgs} node [right] {::msg::ActionExecution(?)} (plan-clients);
      \draw [arrow] (nav-client) -- node [left] {nav2\_msgs::action} node [right] {::NavigateToPose} (nav-server);
      \draw [arrow] (do-work.7) -- node [above] {float32 x, y} (nav-client.148);
      \draw [arrow] (nav-client.154) -- node [below] {bool success} (do-work.353);
      \draw [arrow] (get-pose.353) -| (pose-server.95);
      \draw [arrow] (pose-server.85) |- node [above left] {float32 x, y} (get-pose.7);

      \draw [arrow] (do-work) -- (get-pose);
      \draw [arrow] (get-pose) -- (continue-work);

  \end{tikzpicture}
  \caption[Nodes interaction deepening]{\textbf{DO WORK}, \textbf{GET POSE} and \textbf{CONTINUE WORK} functions are shown in the next page.}
\end{figure}

\newpage

\subsection*{DO WORK function}

\noindent\begin{lstlisting}[language=C, style=CStyle, caption={Starts navigation with the given goal coordinates}]
if(!is_initial_distance_set_){
  room_goal_ = current_arguments_[GOAL_INDEX];
  goal_position_ = labels_.at(room_goal_).at("center").get<std::vector<double>>();
  
  auto request = std::make_shared<planning_bridge_msgs::srv::StartNavigation::Request>();
  request->x = goal_position_[0];
  request->y = goal_position_[1];
  
  start_navigation_client_->async_send_request(request,
  [this](rclcpp::Client<planning_bridge_msgs::srv::StartNavigation>::SharedFuture future){
    auto result = future.get();
    if(result->success) {
      RCLCPP_INFO(get_logger(), "Navigation started");
      get_pose();
    } else {
      RCLCPP_ERROR(get_logger(), "Navigation failed");
    }
  });
  
} else {
  get_pose();
}
\end{lstlisting}

\subsection*{GET POSE function}

\noindent\begin{lstlisting}[language=C, style=CStyle, caption={Asks pose and sets initial and current poses}]
current_pose_client_ptr_->async_send_request(
  std::make_unique<planning_bridge_msgs::srv::CurrentPose::Request>(),
  [this](rclcpp::Client<planning_bridge_msgs::srv::CurrentPose>::SharedFuture future) {
    auto result = future.get();
    if(!is_initial_distance_set_){
      initial_distance_ = std::sqrt(std::pow(result->x - goal_position_[0], 2) +
      std::pow(result->y - goal_position_[1], 2));
      is_initial_distance_set_ = true;
      current_distance_ = initial_distance_;
    } else {
      current_distance_ = std::sqrt(std::pow(result->x - goal_position_[0], 2) +
      std::pow(result->y - goal_position_[1], 2));
      if(current_distance_ > initial_distance_){
        float temp = current_distance_;
        current_distance_ = initial_distance_;
        initial_distance_ = temp;
      }
    }
    continue_work();
    });
\end{lstlisting}

\newpage

\subsection*{CONTINUE WORK function}

\noindent\begin{lstlisting}[language=C, style=CStyle, caption={Sends feedback until the goal is reached}]
progress_ = 1.0 - (current_distance_ /
initial_distance_);

if (progress_ < 0.95) {
  send_feedback(progress_, "Move running");
} else {
  finish(true, 1.0, "Move completed");

  progress_ = 0.0;
  is_initial_distance_set_ = false;
  std::cout << std::endl;
}

std::cout << "\r\e[K" << std::flush;
std::cout << "Moving ... [" << std::min(100.0, progress_ * 100.0) << "%]" << std::flush;
  \end{lstlisting}

\end{document}
